\section{Discussion}
\label{sec:discussion}

\subsection{Three shapes, three mechanisms}

The single-friction sweeps map cleanly to three distinct mechanisms by which information or timing structure interferes with learned collusion. Latency delays signals. It destroys reactive coordination at $L=1$ but leaves a residual focal point coordination that does not require timely monitoring (the $\Delta\approx 0.16$ floor for $L\geq 5$). Noise corrupts signal value rather than its timing. At high $\sigma$ the perceived state is nearly independent of the rival's actual price, so policies on that state encode no rival relevant information. Both reactive and focal point mechanisms are destroyed, and there is no residual floor. Async restructures the game. Under strict alternation each agent's effective state visitation rate is halved within a fixed iteration budget, and the algorithm has more difficulty discovering the high price equilibria of the simultaneous baseline.

The fact that the three frictions act on three different parts of the learning machinery is what makes their interaction non-trivial. Latency leaves the value of the rival side observation correct (just stale, useful for focal point coordination), noise leaves the timing correct (so an agent could in principle de-noise across many observations, even if the Q-learner does not actually do so), and asynchrony leaves both timing and value of the per-period observation intact but reduces the effective number of distinct learning episodes. These channels are orthogonal in the mechanical sense that each acts on a separate part of the learning system. But mechanical orthogonality is not the same as effect orthogonality. Two frictions can have non-overlapping mechanisms and still produce a joint $\Delta$ that is not the sum or product of their single-friction effects, and the combined-friction results show that all three pairwise combinations interact in exactly this way.

\subsection{Latency plus noise, compounding}

The combined-friction design's strongest finding is that latency and noise compound super-multiplicatively. $\Delta_{\texttt{F110}} = 0.059$ falls below both the multiplicative ($0.124$) and minimum ($0.210$) predictions. The two frictions act on different dimensions of the rival side observation. Latency acts on its time of arrival, noise on its value. Either alone leaves some learning signal intact. Combined, neither timing nor value is reliable, and the residual focal coordination that survives pure latency is destroyed by the noise component. In economic terms, this means information degradation rules behave as policy complements. Combining a delayed disclosure rule with a noise injection rule destroys more collusion than either rule alone, and more than naive analysis would suggest.

\subsection{Asynchrony as a positive interaction}

The most surprising result is that asynchrony, despite reducing $\Delta$ in isolation, partially restores it when imposed jointly with the other frictions. The mechanism we tentatively offer is that strict alternation imposes a one-period structural commitment on the rival, giving the moving agent a known target against which to optimise. When the rival-side signal is itself corrupted or stale, this structural commitment partially substitutes for the missing information. \citet{brown2023} provide empirical support for the same logic outside the simulation setting. In online retail data, firms running pricing algorithms with asymmetric update frequencies sustain higher equilibrium prices than firms updating in lockstep, consistent with the idea that timing structure can substitute for explicit coordination. A competing learning stability explanation cannot be ruled out at this point, and the off-path deviation tests in Section~\ref{sec:factorial-offpath} were designed to falsify either story.

\subsection{A simple model of the async-rescue mechanism}
\label{sec:async-rescue-mechanism}

The empirical pattern (large positive $\beta_{LT}$, smaller positive $\beta_{\sigma T}$) is striking enough that it merits a more explicit mechanism argument rather than interpretation alone. The argument that follows is not a formal theorem but a logical decomposition that derives the qualitative pattern from the structure of the friction itself.

\paragraph{Setup.} Define a rival's update event as a calendar period in which the rival's price actually changes, that is, a period in which the rival is the mover under the asynchrony schedule. Under simultaneous play ($T=0$) every calendar period is an update event for the rival. Under $T$-period strict alternation each agent's update events occur every $2T$ calendar periods (with $T=1$, every other calendar period).

\paragraph{Mechanism claim (informal).} Under observation latency $L$, an agent's perception of the rival's price is stale by $L+1$ calendar periods. In rival update events, the same observation is stale by approximately
\begin{equation*}
    \text{rival-update staleness} \;\approx\; \frac{L+1}{2T} \quad \text{(for $T\geq 1$, taken as 1 when this is less than 1).}
\end{equation*}
At $T=0$ this collapses to $L+1$ rival updates. At $T=1$ it is roughly $(L+1)/2$. The same calendar-period latency therefore corresponds to fewer rival decisions missed when the rival updates less often.

\paragraph{Implication 1 (latency rescue).} Take a Q-learner that observes the rival's price with latency $L$, meaning the value it sees is the rival's price from $L$ periods ago. Under simultaneous play, $L$ calendar periods of staleness correspond to $L$ missed rival decisions, because the rival has been moving every period in between. Under strict alternation, where the rival only changes price every other period, those same $L$ calendar periods now correspond to only $L/2$ missed rival decisions. The observation is just as stale in calendar time, but it represents far fewer actual rival moves that the agent missed. A Q-learner that struggled to coordinate with a fast-moving rival on stale data can therefore do much better when the rival is moving more slowly, even though the observation delay itself has not changed. The effect should be largest when $L$ is large, because that is when the calendar staleness was hurting the most. Empirically, $\beta_{LT} = +0.601$ is by far the largest single interaction in the combined-friction design, exactly as this mechanism predicts.

\paragraph{Implication 2 (noise: weaker, but same direction).} The same logic helps with noise, but more weakly. Asynchrony does not reduce the noise on any single observation. Each one still has standard deviation $\sigma$ around the rival's true price. What changes is that the rival's true price now persists for two consecutive periods rather than just one. So when an agent observes the rival at period $t$ and again at period $t+1$, both readings are noisy estimates of the same underlying price, rather than two noisy estimates of two different prices. The cluster of observations the Q-learner accumulates for a given rival action builds up around the right value more cleanly, which acts like a mild form of denoising by repetition. Because asynchrony does not reduce the noise on any individual observation directly and only helps through this indirect repetition channel, the rescue should be weaker than for latency. $\beta_{\sigma T} = +0.245$ is roughly half the size of $\beta_{LT} = +0.601$, exactly as the mechanism predicts.

\paragraph{Implication 3 (the cost of asynchrony in isolation).} Asynchrony has a real cost when imposed alone. Under strict alternation each agent only moves every other period, so within the same iteration budget each agent gets half as many turns to update its $Q$-table. Learning is slower and the converged strategies are less polished. That cost shows up in the single-friction sweep, where $\Delta$ drops from $0.848$ to $0.378$ at $T=1$. The rescue mechanism therefore involves a trade-off. With async as the only friction, the slow-learning cost dominates and $\Delta$ falls. With latency or noise also present, the benefit of having the rival's price be less out of date (because it now persists longer between updates) outweighs the slow-learning cost, and $\Delta$ ends up higher than the single-friction values would predict.

\paragraph{Why the off-path test is consistent.} If the rescue mechanism is genuinely about the rival's perceived state being less out of date, then a forced deviation should still trigger learnt punishment behaviour in the rescue cells, because the learning is happening on real rival information rather than on a focal point lock-in. This is what Section~\ref{sec:factorial-offpath} finds. Retaliation is present in \texttt{F101} and \texttt{F011} and recovers within 5--15 periods, and in \texttt{F111} on a longer timescale of roughly 60--75 periods consistent with the compounding effect of noise on top of the rival staleness reduction. None of the rescue cells shows the spurious focal point snap back that would be expected if the high $\Delta$ were sustained by mere convergence inertia.

\paragraph{What this does not establish.} The mechanism claim above is an informal argument, not a derivation of the equilibrium $\Delta$ from first principles. A fully formal model would have to specify the Q-learning dynamics and the joint state evolution under combined frictions explicitly and either solve for the stationary distribution analytically (out of scope for tabular Q-learning at this state space size) or characterise it through stochastic approximation arguments. What the argument does establish is that the empirical pattern observed in the combined-friction design and the heatmaps is consistent with a single, simple mechanism rather than a coincidence of several independent effects. Future work formalising this argument into a proper theoretical model is a natural extension.

\subsection{Relation to the existing literature}

The strength of the existing simulation literature is control. \citet{calvano2020} show cleanly that independent Q-learning agents can reach cartel-like outcomes without communication, and \citet{calvano2023} provide an off-path diagnostic for whether those outcomes are supported by punishment. The weakness, for the question asked here, is that the baseline environment is deliberately clean: rival prices are observed without delay or measurement error and both agents move on the same clock. It is therefore an ideal laboratory for proving possibility, but less suited to deciding whether real-market information imperfections dampen or preserve collusion.

The one-friction extensions each answer a narrower question well. \citet{calvano2021} and \citet{zhang2025} show that imperfect monitoring and noise can weaken learned collusion, while \citet{klein2021} shows that alternating price setting can still sustain supracompetitive outcomes in a sequential-pricing environment. Their strength is internal clarity: each isolates one channel. Their weakness is also that isolation. A policymaker cannot infer the effect of a delayed disclosure rule, a noisy public feed, and a pricing-window rule by adding up three separate one-friction results, because Section~\ref{sec:results-factorial} shows that those frictions interact non-linearly.

The empirical papers have the opposite trade-off. \citet{assad2024} and \citet{brown2023} study real markets, so their external validity is stronger than any laboratory simulation. They also provide the best evidence that algorithmic pricing and asynchronous update timing matter outside the model. But observational data cannot easily switch one friction on while holding the others fixed, and it cannot directly observe the counterfactual policy combinations studied here. The contribution of this thesis is to connect those two literatures: it keeps the control of the simulation setting, introduces frictions that correspond to real information and timing constraints, and shows that combinations of those frictions can reverse the conclusion one would draw from single-friction experiments. The result identifies a new interaction mechanism that changes how policy instruments should be evaluated.

\subsection{Do these frictions already exist in real markets?}
\label{sec:discussion-existing-frictions}

You might object at this point that the three frictions tested here are not really new constraints to layer onto a Calvano-style market. They describe the conditions of any real market already. Firms observe their competitors through scraping or distribution channels with measurable latency, work from noisy data feeds, and operate on price-update cadences that vary across firms. If single information frictions individually disrupt learned collusion (and the single-friction sweeps show they do), and if those frictions are already present in any actual market, the worry about algorithmic collusion should be substantially smaller than Calvano's clean baseline suggests.

The combined results partly support and partly reject that objection. They support it in the sense that frictions do lower collusion relative to the clean baseline: moving from $\Delta=0.848$ in \texttt{F000} to $\Delta=0.537$ in \texttt{F111} is a reduction of about $0.31$ on the Nash-to-monopoly profit scale. That is not trivial. But $\Delta=0.537$ is still economically high: the algorithms recover more than half of the profit gap between the competitive Nash benchmark and the joint-monopoly benchmark. It is also far above the latency-plus-noise cell \texttt{F110}, where $\Delta=0.059$. Thus the realistic joint-friction case is not as collusive as the clean Calvano baseline, but it is much more collusive than a one-friction or two-information-friction reading would suggest.

The reason is that real markets do not face one friction at a time. They face something closer to a joint cell such as \texttt{F111}. The async-rescue mechanism means that the very combination of frictions a real market has can preserve learned collusion rather than simply destroy it. Reasoning from any one friction in isolation therefore gives the wrong prediction for what happens when all three are present together.

The empirical record so far is consistent with this reading. \citet{brown2023} document that firms running pricing algorithms with asymmetric update cadences in online retail sustain higher equilibrium prices than firms updating in lockstep. The signature matches the async-rescue cell of the combined-friction design, not the single-friction collapse. So the empirical evidence available does not support a "real-world frictions prevent algorithmic collusion" inference.

The right summary is therefore not "real-world frictions prevent algorithmic collusion." It is that single information frictions disrupt learned collusion when imposed in isolation, but the way real-market frictions combine can support it. This is the central contribution of the thesis: it shows that realism does not enter the model as a simple dampening force. Some realistic frictions weaken collusion, but realistic combinations can also create new coordination channels.

\subsection{Policy implication}

\citet{zhang2025} argues for noise injection as a regulatory instrument against algorithmic collusion. The single-friction noise sweep is broadly consistent with that proposal, since a sufficiently large noise level does drive learned collusion to Nash. The combined-friction results qualify the case in two ways that matter for concrete policy design.

First, noise injection becomes substantially more effective when imposed alongside an observation delay rule. $\Delta_{\texttt{F110}} = 0.059$ versus $\Delta_{\texttt{F010}} = 0.500$. The two information degradation rules are policy complements. A regulator who can impose both should expect more disruption than the additive or multiplicative combination would predict, because the rules amplify each other.

Second, structural timing rules (mandated alternation, mandated staggered pricing windows) are policy substitutes for information degradation rules. Imposing the timing rule undoes part of the information rule's effect on learned collusion. A regulator combining "you must publish prices with bounded noise" with "you must commit to prices for a fixed window" should expect less disruption than from the noise rule alone, because the commitment window rule provides exactly the kind of structural stability that lets agents coordinate around corrupted signals.

To make this concrete, consider three real-world policy scenarios.
\begin{itemize}[topsep=2pt, itemsep=4pt]
    \item \textbf{Mandated reporting delay in petrol retailing.} Germany provides a concrete example of a public fuel-price reporting system. Since 31 August 2013, petrol station operators and firms with price-setting authority have been required to report price changes for Super E5, Super E10, and Diesel in real time to the Market Transparency Unit for Fuels, which then passes the data to consumer information services \citep{bundeskartellamtFuels}. A rule delaying competitor access to that official feed would be the policy analogue of $L\geq 1$. The important caveat is enforcement. If firms can simply buy or build a private real-time database through scraping, mystery shopping, or direct commercial data sharing, the public-feed delay is bypassed. A reporting-delay instrument would therefore need access rules or auditing of authorised data services; otherwise it delays only one information channel, not rival observation itself.
    \item \textbf{Noise injection on published prices.} Zhang style mandated obfuscation of published competitor facing price feeds, in markets where firms scrape each other rather than observing directly. Effective on its own, more effective when combined with the delay above. The political and operational reality of this instrument is admittedly awkward. Regulators are not in the habit of mandating that firms publish noisy or distorted information, and a rule that essentially asks "make your prices a bit harder for competitors to read accurately" sits uncomfortably alongside the usual regulatory preference for transparency. So although the simulation results say this is the second most disruptive single instrument available, implementing it would require a regulator willing to deliberately degrade signal quality in a price-sensitive market, which is a stretch from current practice.
    \item \textbf{Mandated price commitment windows.} Requiring firms to commit to a price for, say, a 24 hour window before changing it. This is the analogue of $T\geq 1$. Our results suggest this disrupts collusion when imposed in isolation, but should not be combined with the previous two rules. Doing so partially undoes their effect by giving algorithms exactly the structural stability they need to coordinate around corrupted signals.
\end{itemize}
The bottom line policy lesson is that combinations of information and timing rules need to be evaluated jointly, not as the sum or product of single-instrument analyses. Combined frictions can still reduce collusion relative to the clean baseline: \texttt{F111} has $\Delta=0.537$, well below the baseline $\Delta=0.848$. But the reduction is much smaller than one would infer from the latency-plus-noise cell alone, where $\Delta=0.059$. A policy designer who treats the instruments as mechanically additive would therefore over-estimate the disruption achieved by a package that also includes a timing rule.
